\documentclass{article}
\usepackage[T1]{fontenc}
\usepackage{lmodern}
\usepackage{graphicx}
\usepackage{amsmath,amssymb}
\usepackage[a4paper,margin=2cm]{geometry}
\usepackage[absolute,overlay]{textpos}
\setlength{\TPHorizModule}{1mm}
\setlength{\TPVertModule}{1mm}
\usepackage[indonesian]{babel}
\usepackage{hyperref}
\setlength{\emergencystretch}{2em}
% unit: Halaman 29

\begin{document}

% === Halaman 29 (llm) ===

\begin{itemize}
    \item Misalkan panjang kunci $m = 10$, maka 10 huruf pertama plainteks dienkripsi dengan kunci $K$, setiap huruf ke-$i$ menggunakan kunci $k_i$.
    \item Untuk 10 karakter berikutnya, kembali menggunakan pola enkripsi yang sama.
\end{itemize}

\noindent Contoh: kunci = \texttt{sony} \quad ($m = 4$)\
Plainteks: \texttt{thisplaintexthissecret}\
Kunci: \texttt{sonysonysonysonysonys}

\vspace{5mm}

Setiap huruf plainteks dienkripsi dengan huruf kunci di bawahnya dengan \textit{cipher} abjad tunggal (misalnya menggunakan \textit{Caesar Cipher}):

\begin{align*}
    (t+s) \pmod{26} \\
    (h + o) \pmod{26} \\
    (i + n) \pmod{26} \\
    \text{dst...}
\end{align*}

\end{document}
